\documentclass[10pt,a4paper]{amsart}
\usepackage{amsmath,amssymb,amsthm,graphicx}
\usepackage{amsmath}
\usepackage{amssymb}
\usepackage{hyperref}
\usepackage[utf8]{inputenc}
\usepackage{csquotes}
\usepackage[foot]{amsaddr}


\usepackage{times}
\usepackage{amsmath,amsthm,amssymb}
\usepackage{array}
\usepackage{bbm}
\usepackage{xcolor}
\usepackage{verbatim}
\usepackage{url}
%\usepackage{authblk}
\pagestyle{plain}


\def\ve{\varepsilon}
\def\vf{\varphi}
\def\vn{\varnothing}

\def\AA{\mathcal{A}}
\def\AAA{\mathbb{A}}
\def\BB{\mathcal{B}}
\def\CC{\mathcal{C}}
\def\CCC{\mathbb{C}}
\def\DDD{\mathbb{D}}
\def\DD{\mathcal{D}}
\def\FF{\mathcal{F}}
\def\HH{\mathcal{H}}
\def\II{\mathcal{I}}
\def\III{\mathbb{I}}
\def\MM{\mathcal{M}}
\def\NN{\mathcal{N}}
\def\NNN{\mathbb{N}}
\def\PPP{\mathbb{P}}
\def\QQQ{\mathbb{Q}}
\def\RR{\mathcal{R}}
\def\RRR{\mathbb{R}}
\def\SS{\mathcal{S}}
\def\SSS{\mathbb{S}}
\def\UU{\mathcal{U}}
\def\VV{\mathcal{V}}
\def\WW{\mathcal{W}}
\def\ZZ{\mathcal{Z}}
\def\ZZZ{\mathbb{Z}}
\def\PP{\mathcal{P}}
\def\TTT{\mathbb{T}}
\def\EEE{\mathbb{E}}


%%%%  jb abbreviations
\def\bK{{\bf K}}
\def\bL{{\bf L}}
\def\bU{{\bf U}}
\def\aut{\mathop{aut}}
\def\betabar{{\bf\beta}}
\def\Uscr{\mathcal U}
\def\myreduct{\upharpoonright}
\def\abar{{\bf a}}
\def\bbar{\bf b}
\def\cbar{\bf c}
\def\lg{\mathop{\rm lg}}
\def\nmdeg{\mathop{\rm nmdeg}}
\def\icl{\mathop{\rm icl}}
% \sidebar{..} puts .. in comment form
\newcommand{\sidebar}[1]{\vskip10pt\noindent
 \hskip.70truein\vrule width2.0pt\hskip.5em
 \vbox{\hsize= 4truein\noindent\footnotesize\relax #1 }\vskip10pt\noindent}



 %%%%

\def\QQQa{\mathbb{Q}\langle \Bar{\alpha},\alpha_0 \rangle}

\def\cc{\mathfrak{c}}

\def\a{\alpha}
\def\b{\beta}
\def\g{\gamma}
\def\d{\delta}
\def\e{\varepsilon}
\def\z{\zeta}
\def\l{\lambda}
\def\k{\kappa}
\def\o{\theta}
\def\r{\rho}
\def\s{\sigma}
\def\t{\tau}
\def\w{\omega}
\def\D{\Delta}

\def\nmdeg{\operatorname{nmdeg}}
\newcommand{\Fmb}[1]{F_{Mb}\left( #1  \right)}
\newcommand{\indf}[1]{F_{ind}\left( #1  \right)}
\newcommand{\fcb}[1]{F_{cb}\left( #1  \right)}
\newcommand{\ficb}[1]{F_{Icb}\left( #1  \right)}





\DeclareMathOperator{\alg}{alg}
\DeclareMathOperator{\RU}{RU}
\DeclareMathOperator{\RM}{RM}
\DeclareMathOperator{\RD}{RD}
\DeclareMathOperator{\tp}{tp}
\DeclareMathOperator{\acl}{acl}
\DeclareMathOperator{\dcl}{dcl}
\DeclareMathOperator{\Cb}{Cb}
\DeclareMathOperator{\dtr}{\delta-trdeg}


\def\id{\operatorname{id}}
\def\Frac{\operatorname{Frac}}
\def\Const{\operatorname{Const}}
\def\spec{\operatorname{Spec}}
\def\span{\operatorname{span}}
\def\exc{\operatorname{Exc}}
\def\Div{\operatorname{Div}}
\def\cl{\operatorname{cl}}
\def\mer{\operatorname{mer}}
\def\trdeg{\operatorname{trdeg}}
\def\ord{\operatorname{ord}}




\newtheorem{theorem}{Theorem}[section]
\newtheorem{thm}[theorem]{Theorem}
\newtheorem{question}[theorem]{Question}
\newtheorem{lemma}[theorem]{Lemma}
\newtheorem{cor}[theorem]{Corollary}
\newtheorem{prop}[theorem]{Proposition}
\newtheorem{fact}[theorem]{Fact}
\newtheorem{claim}[theorem]{Claim}
\newtheorem{subclaim}[theorem]{Subclaim}
\newtheorem{conj}[theorem]{Conjecture}
\newtheorem{cond}[theorem]{Condition}
\newtheorem{observation}[theorem]{Observation}
\newtheorem{notation}[theorem]{Notation}

\theoremstyle{definition}
\newtheorem{definition}[theorem]{Definition}
\newtheorem{example}[theorem]{Example}
\newtheorem{xca}[theorem]{Exercise}
\newtheorem{ques}[theorem]{Question}
\newtheorem{note}[theorem]{Note}

\theoremstyle{remark}
\newtheorem{remark}[theorem]{Remark}
\newcommand{\m}{\mathbb }
\newcommand{\mc}{\mathcal }
\newcommand{\mf}{\mathfrak }
\newcommand{\mb}{F_{Mb} }
\newcommand{\G}{\mathrm{G} }
\newcommand{\U}{\mathrm{U} }



\def \ta {\tau_{\mathcal{D}/\Delta}}
\def \D {\Delta}
\def \DD {\mathcal D}

\newcommand{\gen}[1]{\left\langle#1\right\rangle}
%\renewcommand{\restriction}{\mathord{\upharpoonright}}

\newcommand{\pd}[2]{\frac{\partial #1}{\partial #2}}
\newcommand{\td}{\text{tr.deg.}}
\newcommand{\pp}{\partial }
\newcommand{\pdtwo}[2]{\frac{\partial^2 #1}{\partial #2^2}}
\newcommand{\od}[2]{\frac{d #1}{d #2}}
\def\Ind{\setbox0=\hbox{$x$}\kern\wd0\hbox to 0pt{\hss$\mid$\hss} \lower.9\ht0\hbox to 0pt{\hss$\smile$\hss}\kern\wd0}
\def\Notind{\setbox0=\hbox{$x$}\kern\wd0\hbox to 0pt{\mathchardef \nn=12854\hss$\nn$\kern1.4\wd0\hss}\hbox to 0pt{\hss$\mid$\hss}\lower.9\ht0 \hbox to 0pt{\hss$\smile$\hss}\kern\wd0}
\def\ind{\mathop{\mathpalette\Ind{}}}
\def\nind{\mathop{\mathpalette\Notind{}}}
\numberwithin{equation}{section}
\newcommand{\is}{^{p^ {-\infty}}}
\newcommand{\QQ}{\mathbb Q}
\newcommand{\fh}{\mathfrak h}
\newcommand{\SL}{\operatorname{SL}}


\title{Simple Homogeneous Structures and Indiscernible Sequence Invariants}
\author[1]{John T. Baldwin}
\author[2]{James Freitag}
\author[3]{Scott Mutchnik}
%\affil[1]{Department of Mathematics, Statistics, and Computer Science, University of Illinois at Chicago}
%%%\affil[2]{Department of Biology, University Y}
%%\author{ John T. Baldwin}
%\author{James Freitag} 
\address{Department of Mathematics, Statistics, and Computer Science, University of Illinois at Chicago }
\email{jbaldwin@uic.edu}
\email{jfreitag@uic.edu}

%\author{Scott Mutchnik}
%\address{Department of Mathematics, Statistics, and Computer Science, University of Illinois at Chicago}
\email{mutchnik@uic.edu}

\thanks{JB was partially supported by Simons travel grant G3535. JF was partially supported by NSF CAREER award number 1945251. SM, who is additionally affiliated with the Institut de mathématiques de Jussieu – Paris Rive Gauche, was supported by the NSF under Grant No. DMS-2303034, and by the European Union’s Horizon 2020 research and innovation programme under the Marie Sklodowska-Curie grant agreement No 101034255.}

\date{\today}
\title{An omega-categorical supersimple theory whose fibre dimension strictly exceeds the non-minimality degree}
\author{John Baldwin, James Freitag, Scott Mutchnik}
\date{arxiv:2405.08211v3; CC BY 4.0}
\begin{document}
\maketitle
\noindent\textit{Editorial note.} All delivered mathematical text is the author's own, copied from the pinned publisher source sources/logic-fl03-simple-homogeneous-structures-and/pinned-original-source.tex; no statement or proof step was rewritten, reordered or removed.\par\medskip
\noindent\textit{Editorial note.} Preamble provenance: the retained preamble is the source's own layout adapter, if any, plus the author's macro block at lines 3-205 of sources/logic-fl03-simple-homogeneous-structures-and/pinned-original-source.tex (everything between the publisher class metadata and \textbackslash{}begin\{document\}, comments excluded). Nothing was added, no macro was invented, and the publisher document class is deliberately not reproduced; the excerpt serves only as preamble, never as a span coordinate.\par\medskip
\noindent\textit{Editorial note.} Metadata provenance: source\_id and source\_version were taken from sources/logic-fl03-simple-homogeneous-structures-and/license-observation.json; the title from SOURCE.json; the author list from the source metadata record; the source URL from the first non-empty of source\_url, source\_urls[0] and the evidence's fixed\_version\_url/license\_url. The licence evidence file used by this card is sources/logic-fl03-simple-homogeneous-structures-and/license-observation.json, recorded in license.evidence\_path so that the binding target is explicit.\par\medskip
\noindent\textit{Editorial note.} External references are navigation pointers only; every deliberately kept step of the author's argument stays inside the counted spans. The displayed numbers of those pointers are the author's own numbering under the author's \textbackslash{}newtheorem declarations (section.index as declared), computed mechanically by the lane's scanner.\par\medskip
\noindent\textit{Editorial note.} Bibliography provenance: as for the other cards of this source, the source directory stages no bibliography at all, so the single citation in the counted proof is resolved with the bibtex database of the ACL lane's copy of the same arXiv version (focused-seeds/ACL/sources/simple-homogeneous-2405.08211v3/refs.bib, source\_id arxiv:2405.08211v3, source\_version 2405.08211v3); the key Ev02 was verified present (line 1780 of that file) and no entry was edited.\par\medskip
\noindent\textit{Editorial note.} The counted statement and proof contain the author's own commented-out draft passages, which are LaTeX comments and carry no typeset content; they are part of the staged source text.\par\medskip
\noindent\textit{Allowed standard prerequisites (editorial).}\begin{itemize}
\item Evans's construction of *-generic structures from a class C\_f with predimension, and the criteria for the generic to be well-defined and simple. Used inside the counted proof at lines 1758-1766; the individual amalgamation check follows Remark 3.8 of the cited Evans paper. Produces the theory S\_n and justifies that the generic exists and is simple even though the criteria hold only for x >= n+2.
\item The relation between SU-rank, non-minimality degree and predimension in this setting. Used inside the counted proof at lines 1769-1771 and 1778-1781. Turns the rank computation and the predimension estimates into the two bounds nmdeg(p) <= n and nmdeg(p) >= n.
\end{itemize}\medskip
\begin{theorem}\label{simpfind}
Let $n < \omega$. There is an $\omega$-categorical supersimple theory $S_n$
admitting a type $p$ with $\mathrm{nmdeg}(p) = n$ and $F_{ind}(p) >
\mathrm{nmdeg}(p)$.

\end{theorem}

\begin{proof} 
Let $\mathcal{L} = \{R\}$ for $R$ an
$k$-ary relation with $k = n+1$, and $\delta(A) = 2|A| - \frac{|R(A)|}{k!}$. Then let $S_n$ be the theory of the *-generic $M$ for the class $\mathcal{C}_f$ where $f(x) = 2x -2$ for $x \leq
n+2$, and $f(x) = 2(n+2) -2 + (\frac{\mathrm{log}(x)}{\mathrm{log}(3)}-
\frac{\mathrm{log}(n+2)}{\mathrm{log}(3)})$ for $x \geq n + 2$.


Now the criteria for simplicity and well-definedness of the $*$-generic of $\mathcal{C}_f$, $0 < f'(x) < \frac{1}{x}$ and $f(3x)
\leq f(x) + 1$, will hold for all $x \geq n+2$, but not for all $x \geq 0$.
However the $*$-generic of $\mathcal{C}_f$ is still well-defined and simple. As in Remark 3.8 of \cite{Ev02}, we may
individually check the instances of amalgamation where some of the structures
being amalgamated have size at most $n+1$, since these structures either have
no relations, or have size $n+1$ and one relation. Moreover, because $f(x) \geq 2$ for $x \geq k$, we again
see that there is only one type $p \in S^{1}(\emptyset)$. And because $f(x) \leq 2x -1$ for $x \leq k$, so there is again a structure in $\mathcal{C}_{f}$ of size $k$ with exactly one instance of $R$ up to permutation, $\mathrm{SU}(p) = 2$ with
$\mathrm{nmdeg}(p) \leq k-1$, so here $\mathrm{nmdeg}(p) \leq n+1-1 = n$.

To see  $\mathrm{nmdeg}(p) \geq n$, so $\mathrm{nmdeg}(p) = n$, we show there
is   a unique type in $S^{n}(\emptyset)$ realized by distinct $a_{1}, \ldots,
a_{n}$.
%  We
%also see that, even though it is no longer the case that $f(k) > 2n$,
Because $f(x) > 2n \geq d_{0}(a_{1}, \ldots,
a_{n})$ for $x >n + 1$
$|\mathrm{acl}(a_{1}, \ldots, a_{n})| \leq n+1$. But $|\mathrm{acl}(a_{1} \ldots a_{n})| = n+1$ is
impossible, because the least predimension of a set of size $n+1$ is $2n+1$, which is larger than $d_{0}(a_{1} \ldots a_{n}) = 2n$.

%\sidebar{Why?because  the least predimension of a set of size $n+1$ is $2n +
%1
%> d_{0}(a_{1} \ldots a_{n})$.
% So $a_{1} \ldots a_{n}$ is closed. Thus, $\nmdeg(p) \geq n$ and  $\nmdeg(p) =n$.}



%\sidebar{ likely false start apr 23.
%  Now to show
%that $n=\mathrm{nmdeg}(p) < F_{ind}(p)$, we show the only indiscernible
%sequence $\langle a_{1}: i\in \omega\rangle$ in $p$ is the  infinite
%independent sequence of unrelated points. If $\neg R(a_{1} \ldots a_{n+1})$,
%$d_{0}(a_{1} \ldots a_{n+1}) = 2n + 2$.
%%Suppose $a_{n+1} \nind a_{1},
%%\ldots a_{n}$.
%Since $f(x) > 2n + 2$ for $x > n+2$, $|\mathrm{acl}(a_{1} \ldots a_{n+1})|
%\leq n+2$, so either $a_{1}, \ldots a_{n+1}$ is closed or
%$|\mathrm{acl}(a_{1} \ldots a_{n+1})| = n+2$ and
%$d_{0}(\mathrm{acl}(a_{1}\ldots a_{n+1})) \geq f(n+2) = 2n + 2$.}




  Now to show
that $n=\mathrm{nmdeg}(p) < F_{ind}(p)$, note that if $ \{a_{i}\}_{i < \omega}$ is a nonconstant indiscernible sequence, either every
$n+1$-tuple satisfies $R$ or none does. The first is clearly impossible since
long enough such sequences have negative predimension. In the second case,
$d_{0}(a_{1} \ldots a_{n+1}) = 2n + 2$. Because $f(x) > 2n + 2$ for $x >
n+2$, $|\mathrm{acl}(a_{1} \ldots a_{n+1})| \leq n+2$, so either $a_{1},
\ldots a_{n+1}$ is closed or $|\mathrm{acl}(a_{1} \ldots a_{n+1})| = n+2$ and
$d_{0}(\mathrm{acl}(a_{1} \ldots a_{n+1})) \geq f(n+2) = 2n + 2$. For either
option $a_{n+1} \ind a_{1}, \ldots a_{n}$ and $F_{ind}(p) > n$.
\end{proof}

\bibliographystyle{plain}
\bibliography{refs}
\end{document}
